%! Author = lili
%! Date = 4/27/26

\chapter{Tutorium 20.04.2026}\label{ch:tutorium-20.04.2026}
\untit{Präsenzaufgaben 1.1: Wahrscheinlichkeitsräume \hfill 20.--24.\ April 2026}

\section{PÜ 1.1.4 - Faire Münze 3-mal Werfen}\label{sec:pu-1.1.4}
Eine faire Münze wird dreimal geworfen.

\begin{auf}[Geben Sie einen geeigneten Ereignisraum $\eraum$ für dieses Experiment an, der
es erlaubt, die folgenden Fragen zu beantworten.]
    NOCH KEINE LÖSUNG % TODO NOCH KEINE LÖSUNG
\end{auf}

\begin{auf}[Geben Sie eine für dieses Zufallsexperiment geeignete
Wahrscheinlichkeitsfunktion $\wfk$ an. Begründen Sie Ihre Wahl.]
    NOCH KEINE LÖSUNG % TODO NOCH KEINE LÖSUNG
\end{auf}

\begin{auf}[Definieren Sie Ereignisse $A$ und $B$ als Teilmengen von $\eraum$ derart, dass $A$ das Ereignis ist, dass
die Münze beim zweiten Wurf Kopf zeigt; $B$ das Ereignis ist, dass die Münze genau zweimal Kopf zeigt.]
    NOCH KEINE LÖSUNG % TODO NOCH KEINE LÖSUNG
\end{auf}

\begin{auf}[Berechnen Sie mit Hilfe der Wahrscheinlichkeitsfunktion $\wfk$ die
Wahrscheinlichkeiten von $A$ und $B$ sowie vom Schnitt $A \cap B$ und von der
Vereinigung $A \cup B$.]
NOCH KEINE LÖSUNG % TODO NOCH KEINE LÖSUNG
\end{auf}


\section{PÜ 1.1.5 - Faire Münze n-mal Werfen (Zusatzaufgabe für die Flinken)}\label{sec:pu-1.1.5}
Eine faire Münze wird $n$-mal geworfen, wobei $n > 2$ beliebig sei
\begin{auf}[ Beantworten Sie die Fragen (a)--(d) aus Aufgabe~1.1.4 auch für das $n$-mal
wiederholte Werfen der fairen Münze.]
    NOCH KEINE LÖSUNG % TODO NOCH KEINE LÖSUNG
\end{auf}

\begin{auf}[Berechnen Sie die Wahrscheinlichkeit, dass die Münze mindestens einmal Kopf
zeigt.]
    NOCH KEINE LÖSUNG % TODO NOCH KEINE LÖSUNG
\end{auf}


\section{PÜ 1.1.6 - Mensch ärgere dich nicht (Zusatzaufgabe für die Flinken)}\label{sec:pu-1.1.6}
Beim Spiel \emph{Mensch ärgere dich nicht} würfelt der Spieler zunächst einmal.
Dann wird die Figur des Spielers auf dem Brett um die entsprechende Anzahl an
Feldern vorgerückt. Falls der Spieler eine $6$ gewürfelt hat, so darf er anschließend
erneut würfeln. Wieder wird die Figur um die entsprechende Anzahl an Feldern
vorgerückt. Dieser Vorgang wird wiederholt, bis der Würfel zum ersten Mal keine $6$
zeigt. Uns interessiert nur, um wie viele Felder die Figur insgesamt vorgerückt wird.

\begin{auf}[Geben Sie einen geeigneten Ereignisraum $\eraum$ für dieses Zufallsexperiment an.
Finden Sie weitere sinnvolle Wahlen für $\eraum$?]
    NOCH KEINE LÖSUNG % TODO NOCH KEINE LÖSUNG
\end{auf}